\section{Round-twenty-nine reconstruction: universal coarea anchors and exact coefficient ratios}
\label{sec:r29-b1}

A high-probability good-block statement cannot by itself produce an
integrable Fourier tail on its exceptional event.  The present proof separates
two mechanisms.  Predictable good blocks provide central and compact-frequency
contraction.  A deterministic finite coarea-anchor cover provides polynomial
decay in every configuration, including every exceptional component.

\subsection{Nested conditioning and probability normalization}

Let $P^0_{\varepsilon,N}$ be the normalized exact-number hard-sphere
preparation and let $G_N=(G_{N,Z},G_{N,R})$ be the quotient constraint, with
lattice component in $E_Z$ and continuous component in $E_R$.  For a path
source $H$ and centred multiplier $\theta$ put
\[
 Z_{\varepsilon,N}(H,\theta)
 =E_{P^0_{\varepsilon,N}}
   \exp\{\mu_\varepsilon H+\theta\cdot(G_N-Ng_0)\},
 \qquad Z_{\varepsilon,N}(0,0)=1.                        \tag{B1.1}
\]
Particle labels are divided into deterministic blocks $B_1,\ldots,B_J$.
They are exposed in the fixed order
\[
 \mathcal F_j=\sigma(B_1,\ldots,B_j),
 \qquad \mathcal F_0\subset\cdots\subset\mathcal F_J.    \tag{B1.2}
\]
All conditional products below use this increasing filtration; no block is
conditioned on its full complement and then reused in a nonnested tower.

The B2 grand-canonical pressure $P_{\rm GC}(s,H)$ is evaluated before exact
coefficient extraction.  The source-dependent saddle $s_H$ solves
\[
 \nabla_sP_{\rm GC}(s_H,H)=g_0.                           \tag{B1.3}
\]
Numerator and denominator use their own saddles $s_H$ and $s_0$.

\subsection{Predictable good blocks}

A block is good if, conditionally on $\mathcal F_{j-1}$, its quotient
increment has covariance at least $c_*I$ on the conservation quotient and its
lattice support generates the prescribed lattice class.  The regular B2
preparation includes a packing condition on boxes of side
$\ell_\varepsilon\gg\varepsilon$:
\[
 \sup_Q\#\{\hbox{exterior hard-core faces meeting }Q\}\le C_m           \tag{B1.4}
\]
for every block box $Q$ and fixed block size $m$.  This is an explicit
preparation hypothesis and is stable under the B2 regular recovery.

\begin{lemma}[Predictable good-block lower bound]
\label{lem:r29-b1-good}
There are $p_*,\theta,c>0$ such that, until $\theta N$ blocks have been
exposed,
\[
 P\{B_j\hbox{ is good}\mid\mathcal F_{j-1}\}\ge p_*      \tag{B1.5}
\]
and
\[
 P\{\#\hbox{good blocks}<p_*\theta N/2\}\le e^{-cN}.      \tag{B1.6}
\]
The constants are uniform on a compact source chart.
\end{lemma}

\begin{proof}
Condition (B1.4) gives finitely many signed-distance charts for the current
block.  The B2 cluster estimate compares the conditional density with the
Maxwellian block density on the allowed region.  A finite family of legal
block configurations spans the quotient constraint with positive covariance
and has a uniform clearance margin.  Their neighbourhoods retain positive
conditional mass, proving (B1.5).  The predictable Bernoulli exponential
supermartingale gives (B1.6).
\end{proof}

\subsection{A universal coarea anchor on every component}

Reserve a deterministic set $A_N$ of $m_0$ initial velocity/position variables
whose law has a $C^{s+2}$ density after the exact conservation quotient.  Let
$F_N$ be the continuous constraint map from these variables and the exposed
exterior to $E_R$.  The \emph{universal anchor certificate} is a finite family
of minors $M_1,\ldots,M_L$ and a boundary-flat partition
$\chi_1+\cdots+\chi_L=1$ such that
\[
 \chi_\ell\ne0\quad\Longrightarrow\quad |M_\ell|\ge c_A>0              \tag{B1.7}
\]
for every regular exterior and every anchor configuration in the support.
The singular hard-core faces are included as separate coarea charts, not
removed probabilistically.  Such a family is obtained by quotienting exact
conservation directions and selecting, at each point, a maximal minor of the
finite-dimensional derivative $DF_N$; compactness of the normalized anchor
support gives a finite subcover.

\begin{lemma}[Configurationwise anchor density]
\label{lem:r29-b1-anchor}
Conditionally on all non-anchor variables, the push-forward of the anchor law
by $F_N$ has a zero-extended $W^{s,1}(E_R)$ density $r_{N,\mathrm{ext}}$ with
\[
 \norm{r_{N,\mathrm{ext}}}_{W^{s,1}}
 \le C_s                                                   \tag{B1.8}
\]
uniformly in $N$, $\varepsilon$, the exterior, and two centred source
derivatives.  Consequently
\[
 |\widehat r_{N,\mathrm{ext}}(v)|
 \le C_s(1+|v|)^{-s}                                      \tag{B1.9}
\]
on every event, including the event in (B1.6).
\end{lemma}

\begin{proof}
On the support of $\chi_\ell$, use the selected minor as coarea coordinates.
The Jacobian is bounded below by (B1.7), while all amplitude derivatives are
bounded by the regular preparation.  Boundary flatness through order $s-1$
removes boundary distributions when the density is extended by zero.
Hard-core boundary charts use tangential coarea coordinates and the flux
factor.  Summing the finite partition proves (B1.8).  Repeated Fourier
integration by parts gives (B1.9).  Centred source derivatives act on the
amplitude and subtract their conditional means; the fixed derivative budget
keeps the same bound.
\end{proof}

\subsection{Full-frequency damping}

Let $\Phi_N(u,v)$ be the characteristic function of the centred quotient
constraint, $u$ in the lattice torus and $v\in E_R^*$.  Good blocks are
integrated successively using (B1.2); the anchor variables are integrated last.

\begin{theorem}[Exceptional-component smoothing]
\label{thm:r29-b1-damping}
For $s>\dim E_R+5$,
\[
 |\Phi_N(u,v)|
 \le C(1+|v|)^{-s}
 \left[e^{-cN\min\{|(u,v)|^2,1\}}+e^{-cN}\right].         \tag{B1.10}
\]
The estimate is uniform under two centred source derivatives.  In particular,
the exceptional term is integrable over the full continuous Fourier space
without assuming that an exceptional configuration contains a good block.
\end{theorem}

\begin{proof}
On the event with linearly many good blocks, the conditional covariance and
aperiodicity estimates multiply along the increasing filtration and give the
Gaussian/annular factor in brackets.  On the complement, use only its
probability $e^{-cN}$.  In both cases Lemma~\ref{lem:r29-b1-anchor} is applied
after the event has been resolved and supplies the factor
$(1+|v|)^{-s}$.  Thus smoothing is configurationwise, not borrowed from a
possibly absent good block.  Centred source derivatives mark at most two
conditional factors and the anchor amplitude; their $N$-scale contributions
cancel against differentiation of the pressure and are bounded by the same
right side.
\end{proof}

\subsection{Mixed local coefficients}

Let $d_Z=\dim E_Z$ and $d_R=\dim E_R$.  At the saddle $s_H$, let
$\Sigma_H$ be the positive covariance on the conservation quotient.  Its
lattice block has covolume $\operatorname{covol}(\Lambda_H)$.

\begin{theorem}[Source-uniform exact-number local theorem]
\label{thm:r29-b1-llt}
For central lattice/continuous deviations $z_N=O(\sqrt N)$,
\[
 P_H\{G_{N,Z}=k_N,\ G_{N,R}\in dy_N\}
 ={\operatorname{covol}(\Lambda_H)
   e^{-\frac1{2N}z_N^T\Sigma_H^{-1}z_N}
  \over(2\pi N)^{(d_Z+d_R)/2}\sqrt{\det\Sigma_H}}
 \left(1+O(N^{-1/2})\right)dy_N.                          \tag{B1.11}
\]
The estimate and two source derivatives are uniform.  Conditioning on a
regular continuous constraint is the corresponding coarea disintegration.
\end{theorem}

\begin{proof}
The central ball uses the third-order expansion of the B2 pressure at $s_H$.
The compact annulus is controlled by the good-block aperiodicity gap.  The
unbounded continuous region is absolutely integrable by
Theorem~\ref{thm:r29-b1-damping}.  Fourier inversion gives (B1.11), and
completion of the Gaussian square gives every conditional Schur covariance.
The coarea formula applies to the resulting continuous density, rather than to
a fixed-width window.
\end{proof}

\subsection{Separate saddles without undefined indices}

Let $A_N(H)$ and $A_N(0)$ be coefficients of the same fixed-dimensional
lattice/coarea constraint.  Let $\mathcal H_H=D_s^2P_{\rm GC}(s_H,H)$ and let
$c_H$ denote the local lattice/coarea amplitude.  Positivity of the
grand-canonical coefficients and the lattice aperiodicity certificate imply a
strict modulus gap away from the real saddle on the selected activity torus.

\begin{theorem}[Dimension-consistent coefficient ratio]
\label{thm:r29-b1-ratio}
Uniformly on a compact source chart,
\[
 {A_N(H)\over A_N(0)}
 =\exp\{N[q(H)-q(0)]\}
   \left({\det\mathcal H_0\over\det\mathcal H_H}\right)^{1/2}
   {c_H\over c_0}\left(1+O(N^{-1/2})\right).              \tag{B1.12}
\]
The common power $N^{-(d_Z+d_R)/2}$ cancels.  No undefined
``$\operatorname{ind}H$'' occurs, and a source does not change the dimension
of a nondegenerate fixed constraint.
\end{theorem}

\begin{proof}
Deform numerator and denominator separately through $s_H$ and $s_0$.  The
positive-coefficient modulus inequality is strict outside the aperiodic real
class, so all other activity arcs are exponentially smaller.  Apply
Theorem~\ref{thm:r29-b1-llt} and the same coarea map to each coefficient.
Their equal Gaussian dimensions cancel, leaving the Hessian and amplitude
ratio in (B1.12).
\end{proof}

The exact-number preparation now has nested conditioning, a smoothing anchor
on every exceptional component, explicit source scaling, and a correctly
dimensioned separate-saddle ratio.
